\chapter{Program Overview}

This manual is an 8-week applied research program for undergraduate interns working in the Project WiCCED Salt project space. The central training goal is practical: students learn how to turn water-quality observations, chemical screening evidence, public database records, and machine-learning models into reproducible research products that can be inspected, challenged, improved, and handed off.

The curriculum is designed for co-advising by Dr. Chase Stratton and Dr. Gulnihal Ozbay. It connects uafR, cheminformatics, reproducible R workflows, water-quality field and lab interpretation, aquaculture and coastal ecosystem context, and machine learning. It does not treat machine learning as a decoration. Models are introduced only after students understand the variables, sampling context, missingness, QA/QC flags, and evidence limits.

\begin{programbox}{Program Promise}
By the end of eight weeks, each student should have a complete water-quality research package: a feasible research question, a reproducible R project folder, a cleaned water-quality dataset or case dataset, documented QA/QC decisions, optional uafR chemical-screening evidence, an analysis-ready feature table, at least one defensible baseline machine-learning model, model evaluation outputs, source-backed figures and tables, a manuscript-style report, a presentation, and a continuation note.
\end{programbox}

\section{Audience and Prerequisites}

The manual assumes students are beginning undergraduate researchers. Prior experience in R, chemistry, water quality, or machine learning is helpful but not required. The program does assume sustained effort, careful documentation, and willingness to revise work after evidence checks.

Minimum technical prerequisites are:
\begin{itemize}
  \item R and RStudio installed and working.
  \item Access to the uafR package from the provided student bundle or local repository.
  \item Ability to open CSV files, create folders, run R scripts, and read console messages carefully.
  \item Access to either real project data approved for training use or the teaching data included with this manual.
\end{itemize}

Minimum research prerequisites are:
\begin{itemize}
  \item A research log that records decisions, file names, dates, errors, corrections, and interpretation limits.
  \item A clear distinction between observations, processed variables, model outputs, and scientific claims.
  \item A willingness to say ``not supported by the current evidence'' when the data do not justify a claim.
\end{itemize}

\section{How to Use This Manual}

The manual is organized as a working program, not a reference book. The 56-day spine should be followed in order unless the co-advising team deliberately adjusts pacing for a specific project. Weekdays are full internship days. Weekend days are lighter required asynchronous work: reading, annotation, reruns, research-log updates, peer review, project maintenance, or revision.

The curriculum has four simultaneous tracks:
\begin{itemize}
  \item \textbf{Water-quality science:} salinity, conductivity, dissolved oxygen, nutrients, pH, temperature, turbidity, chlorophyll, microbial indicators, contaminants, field metadata, QA/QC, and ecological interpretation.
  \item \textbf{uafR and chemical evidence:} tentative GC-MS annotation, public database enrichment, source-backed traits, evidence trails, and chemical grouping.
  \item \textbf{Machine learning:} feature engineering, train/test design, baseline models, classification, regression, model diagnostics, uncertainty, interpretability, and leakage control.
  \item \textbf{Research communication:} figures, tables, claim-evidence-limitation writing, reproducibility statements, presentations, and handoff notes.
\end{itemize}

\section{What Students Produce}

Every student product should be saved inside a reproducible project folder. The final folder should contain:
\begin{itemize}
  \item Project brief and question-scope record.
  \item Raw-data inventory and data-use note.
  \item Cleaned water-quality dataset or approved case dataset.
  \item Water-quality variable dictionary.
  \item QA/QC flag log.
  \item Optional uafR chemical-screening result object and evidence map.
  \item Feature table used for modeling.
  \item Train/test split record.
  \item Machine-learning model card.
  \item Model metrics and diagnostic plots.
  \item At least three publication-style figures or tables.
  \item Manuscript-style report.
  \item Presentation or poster.
  \item Daily research log and weekly review notes.
  \item Final package audit and continuation note.
\end{itemize}

\section{What Mentors Review}

The review process is evidence-centered. Mentors review whether the student can explain what each file is, why each processing step occurred, what evidence supports each variable, what the model can and cannot claim, and whether the final package can be rerun.

Review checkpoints are built into every week:
\begin{itemize}
  \item \textbf{Week 1:} research question, project scope, and risk note.
  \item \textbf{Week 2:} reproducible folder and clean-session rerun.
  \item \textbf{Week 3:} water-quality data dictionary and QA/QC plan.
  \item \textbf{Week 4:} uafR evidence map or chemical-screening decision note.
  \item \textbf{Week 5:} first defensible baseline model and validation plan.
  \item \textbf{Week 6:} feature table, leakage audit, and model comparison.
  \item \textbf{Week 7:} evidence-based interpretation and figure package.
  \item \textbf{Week 8:} final report, presentation, archive, and continuation note.
\end{itemize}

\section{Scientific Integrity Rules}

\begin{cautionbox}{Claims That Must Stay Conservative}
This program trains students to interpret evidence, not to decorate uncertain results. A tentative GC-MS match is not a confirmed chemical identity. A PubChem or KEGG record is not evidence that a compound is present in a water sample. A pathway annotation is not evidence of pathway activity. A hazard label is not a site-specific risk assessment. A machine-learning model association is not proof of causation. A high-accuracy model can still be invalid if it leaks site, time, treatment, or post-outcome information.
\end{cautionbox}

Acceptable wording is specific and limited. For example:
\begin{itemize}
  \item Stronger: ``The tentative GC-MS feature matched caffeine in the library search and uafR linked the query to PubChem CID 2519; this supports follow-up confirmation, not confirmed sample identity.''
  \item Stronger: ``Conductivity, salinity, temperature, and site context were useful predictors in the teaching model under the recorded split; external validation is required before operational use.''
  \item Weaker: ``The site contains caffeine and is biologically active.''
  \item Weaker: ``The model proves saltwater intrusion is causing the observed biological response.''
\end{itemize}

\section{Weekly Time Budget}

A typical weekday contains 6 to 7 active hours:
\begin{itemize}
  \item 30 minutes: check-in, daily objective, research-log setup.
  \item 60 to 90 minutes: concept block or reading discussion.
  \item 90 to 120 minutes: guided technical lab.
  \item 90 to 150 minutes: independent project work.
  \item 45 to 60 minutes: evidence, writing, figure, or documentation block.
  \item 30 minutes: peer or mentor checkpoint.
\end{itemize}

A typical weekend day contains 1 to 3 focused hours. Weekend work should not introduce major new methods unless the student is already ahead. It should consolidate the week: annotate a reading, rerun scripts, revise a worksheet, clean the research log, prepare questions, or review a peer project.

\section{Skill Progression}

The program moves students through four stages:
\begin{itemize}
  \item \textbf{Operator:} runs scripts, checks objects, records errors, and follows data-handling rules.
  \item \textbf{Inspector:} examines missingness, suspicious values, duplicate records, source coverage, and model diagnostics.
  \item \textbf{Analyst:} chooses variables, compares models, creates figures, and explains uncertainty.
  \item \textbf{Researcher:} writes defensible claims, connects evidence to literature, and hands off a reproducible project.
\end{itemize}

\section{Before Week 1}

Before the first day, the project lead should prepare:
\begin{itemize}
  \item The student bundle or local uafR installation path.
  \item This manual as a PDF and source folder.
  \item The \file{training_wicced_water_quality/scripts} folder.
  \item Any approved real water-quality or GC-MS data that students may use.
  \item A decision about whether students will work with real data, teaching data, or both.
  \item A data privacy note for unpublished or partner-owned datasets.
  \item A shared schedule for weekly review checkpoints.
\end{itemize}

Students should arrive with RStudio installed, a folder location selected for project work, and a research log ready to begin. If installation is incomplete, Week 1 can still proceed using paper worksheets, source readings, and the teaching data dictionary.
